\section{Quantum Behaviour from the Deterministic Mesh (P7)}
\label{sec:quantum}

The deepest question for any classical-base ontology is whether it can carry quantum statistics. Bell's theorem (1964) says no local model of independent coins can match the quantum CHSH correlations, unless one gives up the assumption that the measurement settings are statistically independent of the hidden state. A monist mesh with no hidden state denies exactly that assumption, because the settings are made of the same substrate as the hidden state. P7 makes this quantitative.

\subsection{Determinism makes violation possible}

We swept how strongly the measurement settings track the shared hidden state, in a constructed superdeterministic model.

\begin{result}[The cost curve]
As the setting-state correlation rises from independence to full sharing, the CHSH value climbs $1.0, 1.6, 2.3, 3.1, 4.0$, crossing the classical bound of $2$ near half-sharing, passing the quantum Tsirelson bound of $2.83$, and reaching the algebraic maximum of $4$. The price of denying statistical independence is quantified: a given amount of measurement dependence buys a given amount of violation.
\end{result}

\subsection{Aligned bits, not bits}

But the dependence must be of the right kind. We measured the CHSH value alongside the measurement dependence (the mutual information between settings and hidden state, in bits), refining the bound of Hall (2010).

\begin{result}[The currency is aligned bits]
At one bit of measurement dependence, an aligned correlation (tracking the feature the outcomes use) reaches $S = 4$, while a misaligned correlation of the same one bit gives $S = 1$, and a correlation below the measurement resolution gives nothing. So measurement dependence is necessary but radically insufficient. One bit can buy anything from no violation to maximal violation, depending on alignment. The currency of Bell violation is aligned bits, not bits.
\end{result}

\subsection{The residual tension, from dynamics}

The honest frontier is whether a natural mesh produces the aligned correlation without fine-tuning. In a relativistic mesh the shared past of two spacelike measurements (the overlap of their backward light cones) shrinks with their separation.

\begin{result}[Violation decays with separation]
Modelling the shared-past fraction as $e^{-d/\xi}$ in separation $d$, the CHSH value decays from $4.0$ at zero separation to the classical bound within a couple of correlation lengths ($2.58, 1.87, 1.04$ at increasing $d$). Quantum mechanics gives a separation-independent violation. So a natural causal mesh reproduces quantum violation only for nearby measurements, not spacelike ones, unless the correlation is fine-tuned not to decay.
\end{result}

This is the precise residual difficulty of the classical-base reading, now a number rather than a slogan. Vibe Theory makes Bell violation possible from a deterministic mesh, and pins the exact remaining puzzle: why the correlation is aligned and why it would be separation-independent. We made the tension quantitative. We did not resolve it, and we do not claim to. This is the framework's most contested and least settled rung, and it is the one to watch.
